\subsection*{Quels usages de l'IA pour une bibliothèque musicale numérique ? L'exemple de LaDocumenta.eu.}

L'intégration de l' intelligence artificielle au sein de plateformes numériques musicales à l'instar de LaDocumenta.eu, répond à une ambition bien délimitée : l'IA ne doit pas être utilisée pour elle-même dans un souci de conformisme. 

La plateforme LaDocumenta.eu, par exemple, envisage l'utilisation de l'IA comme outil de médiation, de préservation et de transmission du patrimoine musical immatériel de la pratique historiquement informée et de ses \enquote{savoirs tacites\footnote{Cette notion de \enquote{savoir tacite} a été développée dans les candidatures d'appel à financement européen de LaDocumenta.eu}} d'interprétation. Ceux-ci correspondent à des connaissances non verbalisées fondées sur la pratique, l'apprentissage par imitation et le bagage culturel et technique attendu de tout musicien. Ils sont dits \enquote{tacites} car, bien que réputés évidents, ils n'ont jamais fait l'objet d'un apprentissage conscient du musicien. C'est ce type de \enquote{savoirs tacites} que la plateforme LaDocumenta.eu pérenniser.
LaDocumenta.eu est une infrastructure européenne dédiée à la recherche et au patrimoine dont la philosophie est de mettre en parallèle chaque usage métier avec des solutions techniques (IA)- existantes et en développement - dédiées.

Lors du colloque de lancement de la deuxième version de LaDocumenta.eu, organisé en avril 2026 à Prague, les développeurs de la plateforme ont proposé aux collaborateurs de réflechir aux \textit{personae} de LaDocumenta.eu. L'objectif était d'identifier les profils-types des usagers de la plateforme afin de regrouper ces différents profils d'usager par thématiques et répondre au mieux aux besoins. 



Les quatre groupements thématiques des usages de la plateforme sont les suivants : ingénierie du patrimoine, recherche, interprétation et éducation/pédagogie\footnote{Les exemples développés dans les paragraphes suivants sont issus de l'exercice des \textit{personae} proposé aux usagers de LaDocumenta.eu. Les développements sur les solutions d'intelligence artificielle qui pourraient alimenter le fonctionnement de la plateforme sont le fruit des réflexions menées dans le cadre du livrable de scénario d'utilisation de l'IA dans la plateforme LaDocumenta.eu, dans le volume d'annexes de ce mémoire.}. 



Dans le domaine de l'\textbf{ingénierie du patrimoine} recouvrant également l'ingénierie documentaire, les besoins des bibliothécaires et archivistes portent essentiellement sur des tâches de catalogage et de signalement. Dans un contexte de multiplication de la masse documentaire, utiliser l'IA est vu comme une réponse à l'impératif de gain de temps, essentiel pour tout bibliothécaire. 
Ce besoin d'assistant documentaire, apte à gérer des masses de documents, peut être couvert par certains outils d'IA, notamment en ce qui concerne le traitement du langage naturel \textit{NLP : Natural Language Processing}. Ces outils s'envisagent en particulier dans des tâches d'extraction d'entités et de traduction automatique multilingue : l'outil est capable de pré-remplir les formulaires de métadonnées à partir d'un texte de description d'un agent. L'information est alors répartie dans les champs adéquats, ce qui libère un temps considérable à l'agent qui peut alors se concentrer sur des tâches moins répétitives, comme la médiation par exemple. 
Si le recours à ce type d'outils semble simple, il faut toutefois garder à l'esprit que l'implémentation d'un modèle de NLP ne peut s'envisager sans une validation humaine après chaque saisie. Les modèles d'aujourd'hui, s'ils sont déjà compétents pour du langage courant éprouvent, en effet, des difficultés face à un vocabulaire spécialisé - musical par exemple - ou ancien. Par ailleurs, une validation constante des contributions de l'intelligence artificielle permet la correction des erreurs et intrinsèquement l'amélioration du modèle.


Les enjeux de \textbf{recherche} de LaDocumenta.eu concernent le champ spécialisé de la musicologie. Les usagers - chercheurs et musicologues - ont besoin d’étudier des documents complexes tels que des partitions annotées à la main, ou encore des traités d’instruments anciens et de relier ensuite ces documents à d’autres sources dispersées dans les institutions européennes ou mondiales. Ces usages relèvent de la recherche appliquée et sont portés par des collaborations scientifiques solides comme pour les projets REMDM (Répertoire des Écritures Musicales du Département de la Musique)\footnote{REMDM est un projet mené entre 2020 et 2023 à la BnF, avec le soutien de l'IREMUS (Institut de recherche en musicologie) et le concours des laboratoires IRISA (Institut de Recherche en Informatique et Systèmes Aléatoires) de l'université de Lille et L3i de l'université de la Rochelle. L'objectif de ce projet était de constituer \enquote{un répertoire des écritures musicales dues à la fois à des compositeurs et à des copistes identifiés ou anonymes}. \enquote{Parallèlement à la constitution de la base de données, le projet prévoit le développement de méthodologies innovantes d’analyse graphique faisant appel aux technologies les plus avancées en matière de fouille d’images.} La présentation complète de ce projet est consultable à cette adresse : \url{https://remdm.hypotheses.org/le-projet-remdm}.} et Collabscore\footnote{Collabscore est un projet numérisation de partitions musicales soutenu par l'Agence Nationale de la Recherche. Fruit d'un travail collaboratif entre le CNAM (Centre National des Arts et Métiers), la BnF, l'IREMUS (Institut de recherche en musicologie), la fondation Royaumont, l'équipe SHADOC (\textit{Systems for Hybrids Analysis of Documents}) du laboratoire IRISA (Institut de Recherche en Informatique et Systèmes Aléatoires)de l'université de Rennes et l'équipe Algomus du laboratoire CRIStAL (Centre de Recherche en Informatique, Signal et Automatique de Lille) de l'université de Lille, le projet Collabscore (2020-2025) avait pour objectif de rendre exploitable les fonds musicaux numérisés de la BnF et de la Fondation Royaumont qui n'étaient pas encore interrogeables. \url{https://bnf.hypotheses.org/54322}}, initiés par le CNAM (Centre National des Arts et Métiers) et le département de la Musique de la BnF avec le concours d'institutions et de laboratoires spécialisés. Ces projets utilisent les avancées en OMR (\textit{Optical Music Recognition}) et OCR (\textit{Optical Character Recognition}) pour pouvoir encoder du texte et de la musique notée ainsi que l'alignement sémantique par Wikidata et les ontologies CIDOC-CRM, IFLA-LRM. Ces projets à la pointe de la technologie des humanités numériques, s'ils sont très inspirants pour les institutions, nécessitent toutefois un important investissement humain car la préparation des corpus exige un travail d'annotation chronophage. Des solutions de crowdsourcing sont parfois utilisées, comme cela a été le cas pour le projet CollabScore. 


Dans le domaine de l’\textbf{interprétation}, les chefs, les musiciens et les chanteurs sont à la recherche de clés d’interprétation pour nourrir leur jeu. Ils ont besoin d’indications précises sur le choix de l’instrument, la manière de prononcer une langue étrangère, ou encore sur les dispositions à privilégier pour l’acoustique de la salle de concert. Les réponses qu’ils attendent reposent nécessairement sur la recherche musicologique, détaillée dans le point précédent, et sur la mise en place d’outils opérationnels, d’IA par exemple. Une des idées retenues pour LaDocumenta.eu serait de mettre en place un agent conversationnel de type chatbot qui pourrait guider les musiciens dans leur interprétation. Ce chatbot (MentorIA) spécialisé en musicologie est une adaptation du concept des agents conversationnels des bibliothèques universitaires, au premier rang desquels le chatbot T-REX de l'Université de Calgary, cet assistant conversationnel qui combine un LLM souverain et une structure RAG (\textit{Retrieval-Augmented Generation}).
En adaptant ce modèle au cas de LaDocumenta.eu , on pourrait imaginer un chat bot restreignant ses réponses aux seules données qualifiées de LaDocumenta.eu, qui fournirait des conseils d'interprétation. 


Le dernier domaine ciblé par les \textit{personae} des développeurs concerne les domaines de la \textbf{pédagogie} et de l’\textbf{éducation}. Dans le domaine patrimonial, cela peut s’assimiler à la médiation des contenus de la plateforme. Il s’agit donc d’exploiter les ressources de la plateforme pour les présenter au grand public ou aux musiciens non spécialistes de musique historiquement informée, ainsi qu’aux enseignants voulant présenter  cette pratique à leurs étudiants. Les plateformes numériques peuvent avoir recours à l’IA pour proposer des outils de médiation originaux, utilisant par exemple le scrollytelling nourri par des IA génératives adaptées au niveau de connaissance de l’utilisateur afin de cibler le plus finement possible ses besoins et capter ainsi mieux son attention. Cela peut également s’envisager pour des IA génératives audio afin de modifier le rendu sonore d’un enregistrement pour simuler un changement d’instrument, de caractère ou d’époque. 
Une mise en garde s’impose toutefois : si l’IA générative textuelle est d’ores et déjà accessible, l’IA générative audio implique des coûts très élevés, avec des risques pour les enregistrements retouchés (de dégradation du signal sonore d’origine) et des problématiques juridiques de droits d’auteur qui peuvent s’avérer coûteuses.

Si ces quatre axes d'usage dévoilent les besoins théoriques des usagers types de la plateforme, le choix de la technologie retenue repose sur une réflexion précise des parcours métiers qui doit être menée à partir d'un état de l’art des solutions existantes. 
C’est l’objet du livrable de scénario d’utilisation de l’IA que j’ai été amenée à réaliser lors de mon alternance comme chargée de mission de la plateforme LaDocumenta.eu.


\subsection*{Étude de cas : Le livrable de scénario d'utilisation de l'IA de la plateforme LaDocumenta.eu, contexte et méthodologie de rédaction}


Pour envisager l'intégration de l'intelligence artificielle au sein de LaDocumenta.eu, l'identification des usages métiers et du public seule n'était pas suffisante. Identifier les usages devait s'inscrire dans une réflexion plus globale et stratégique sur le choix des outils d'IA intégrés à la plateforme. Cette liste des usages devait s'articuler entre les besois métiers et du public avec les attendus des programmes européens de financement, notamment en ce qui concerne l'intelligence artificielle. 

Ainsi donc, la rédaction du livrable de scénario d'utilisation de l'IA de la plateforme est née de cette double exigence : d'une part fournir au consortium un outil de travail pour la préparation des évolutions IA de la plateforme à visée prospective, et d'autre part préparer le document de résultats attendu par Europe Créative. 

L'enjeu méthodologique de ce travail était de partir des scenarii concrets - métier et public - d'usages de la plateforme et de trouver des solutions d'IA appliquées - existantes ou en construction - à ces usages,
 et non l'inverse. Ainsi, l'idée n'est pas de reprendre des modèles déjà existants pour les calquer sur les besoins de la plateforme, mais d'analyser ces besoins en amont pour y fournir une réponse sur-mesure..

Pour répondre à ces exigences, ce livrable se structure en deux parties: 
\begin{itemize}
	\item La première partie du livrable s'est attelée, comme nous l'avons vu, à lister les usages de la plateforme. Ceux-ci ont été identifiés lors de l'exercice des \textit{personae} qui a permis d'isoler quatre axes stratégiques d'intervention de l'IA : l'ingénierie documentaire patrimoniale, la recherche, l'aide à l'interprétation et la médiation culturelle (pédagogie, éducation).
	\item La deuxième partie du livrable était vouée, quant à elle, à détailler ces usages en proposant des scenarii réalistes de situations d'utilisation de la plateforme. L'enjeu était d'associer à chaque usage un outil d'IA assitant les usagers - professionnels ou grand public - dans leur utilisation de la plateforme.
	Après un premier travail trop centré sur les scenarii eux-mêmes, nous avons enrichi le livrable d'un tableau synthétique associant chaque situation d'utilisation avec son outil d'IA correspondant et un exemple de projet similaire dans l'état de l'art. Il s'agissait de proposer un regard critique sur les solutions énoncées, pour pouvoir enrichir une liste d'outils et évaluer leur faisabilité en fonction des divers scenarii. Chaque scénario d'usage a donc été confronté aux technologies existantes ou en cours de développement, en adossant chaque proposition à un projet de référence (Isidore\footnote{ISIDORE est un « moteur de recherche permettant de découvrir et de trouver des publications, des données numériques et profils de chercheur.e.s en sciences humaines et sociales (SHS) venant du monde entier. Il permet de rechercher dans le texte intégral de plusieurs millions de documents (articles, thèses et mémoires, rapports, jeux de données, pages Web, notices de bases de données, description de fonds d’archives, etc.), des signalements
		événements (séminaires, colloques, etc.). Référence bibliographique : article collectif sur les 10 ans d'ISIDORE : ISIDORE a 10 ans. Zenodo. \url{https://doi.org/10.5281/zenodo.5699997}}, Annif\footnote{ \enquote{la Bibliothèque nationale de Finlande a développé Annif, un programme d’indexation automatique disponible en open source. Après l’avoir « alimenté » avec un vocabulaire SKOS1 (en l’espèce la General Finnish Ontology, YSO) et les métadonnées disponibles sur le site finna.fi, Annif sait désormais comment assigner des termes d’indexation à de nouveaux documents, ce dans plusieurs langues.}. Référence bibliographique : Suominen (Osma), « Annif, l’indexation automatique à la Bibliothèque nationale de Finlande », Arabesques [En ligne], 94 | 2019, mis en ligne le 15 décembre 2019, consulté le 09 septembre 2026. URL :\url{https://publications-prairial.fr/arabesques/index.php?id=605}}, Collabscore ou REMDM) afin d'en évaluer les avantages et inconvénients.
\end{itemize}

La principale difficulté de ce travail a été de déterminer le degré de précision des solutions envisagées. 
C'est pourquoi nous avons fait le choix de proposer un degré de précision variable pour les solutions évoquées dans le tableau synthétique. Il correspond aux différents états d'avancement desdites solutions. 
La méthodologie originale prévoyait de détailler tous les scenarii d'usages mais, après l'étude des solutions, nous avons pris conscience d'une hétérogénéité concernant la maturité des solutions. 
Dans le domaine de l'ingénierie documentaire et patrimoniales, les outils d'IA était déjà industrialisés, accessibles en open source et donc déjà opérationnels, moyennant une adaptation au contexte spécialisé de LaDocumenta.eu. Dans le domaine de la musicologie computationnelle (OMR notamment) en revanche, les solutions proposées relèvent encore de la recherche fondamentale même si des projets aboutis ont été cités (commme Collabscore). Nous avons donc fait le choix de proposer des exemples de scenarii dont la maturité était inégale, cela supposant d'accepter des données parfois partielles. 

Par ailleurs, la rédaction de ce livrable a souligné un autre enjeu majeur de l'intégration de ces outils d'IA, au niveau de la formation des équipes à l'IA. Comme pour le livrable du manuel de dépôt,  l'accompagnement au changement fait donc ici aussi partie intégrante des enjeux de ce livrable qui entend poser un premier jalon dans ce domaine. 

Dans cette perspective, le présent livrable se veut rassurant avec les contributeurs quant à l'intégration de l'IA, non comme remplaçant mais bien comme assistant logistique des agents. Chargés des tâches répétitives (comme le pré-remplissage des métadonnées par exemple), les outils d'IA permettront aux agents de valoriser leur expertise professionnelle dans des opérations de contrôle et de leur libérer du temps de travail.


En outre, le caractère évolutif des outils d'IA, la nature prédictive des scenarii d'IA et le développement encore en cours de la plateforme LaDocumenta.eu interrogent sur le risque d'obsolescence rapide de ce type de livrable. Nous sommes conscients de cette limite même si ce livrable a été pensé comme un jalon de réflexion à un moment donné, nécessaire pour l'envoi des résultats dans le cadre d'Europe Creative. Ce livrable ne présente donc pas de caractère définitif, ce qui explique pourquoi il n'y a pas d'estimation financière associée mais seulement un curseur classant les prévisions de coûts dans les catégories faible, modéré et élevé.
En somme, malgré les limites intrinsèques à la rédaction de ce livrable, il remplit son rôle de document d'évaluation des solutions existantes à un moment T, en confrontant les besoins de la plateforme avec les outils, opérationnels ou en développement, correspondants. 

Cette confrontation des usages métiers et de la recherche appliquée en IA montre qu'une plateforme numérique culturelle spécialisée peut s'emparer d'outils d'IA, sans céder à la vogue de l'IA, pour réfléchir à ses pratiques futures et à son développement.

Ainsi, l'élaboration de ces scenarii d'usage souligne le rôle prépondérant de la plateforme LaDocumenta.eu et de ses futurs outils d'IA dans la transformation des ressources patrimoniales en données exploitables à l'échelle européenne pour la recherche et l'interprétation de la pratique historiquement informée. 

Si les outils d'IA présentent des solutions fiables et prometteuses, il faut néanmoins prendre en compte les aspects juridiques et éthiques qui en découlent. Ces questions émergentes dans le domaine des industries culturelles et créatives feront l'objet du chapitre suivant. 